\section{Модель и метод расчёта}
\label{sec:methods}

\subsection{Геометрия и каналы возбуждения}
\label{sec:geometry}

Рассматривается одиночная металлическая наночастица (МНЧ) — вытянутый сфероид
золота с полуосями $a=\SI{10}{\nano\meter}$ (малая) и $c=\SI{15}{\nano\meter}$
(большая, вдоль оси симметрии) — и одна коллоидная квантовая точка (КТ)
радиуса $r_{\rm QD}=\SI{2}{\nano\meter}$. Система помещена в однородную
непоглощающую среду с $\varepsilon_m=2.25$.
Размеры частиц и проницаемость среды заимствованы из раздела II.D
работы~\cite{shah2013}, где две такие МНЧ образуют антенну
с межчастичным зазором \SI{6}{\nano\meter}, а КТ расположена посередине.
После удаления одной МНЧ исходное положение КТ соответствует
$R=\SI{18}{\nano\meter}$ и поверхностному зазору $g=\SI{1}{\nano\meter}$.
Здесь исследуется одна МНЧ и дополнительно варьируется $g$.
Эффективные параметры плазмона двухчастичной антенны из этой работы
не переносятся на одиночный сфероид.

Зазор $g$ определён как расстояние между поверхностями КТ и МНЧ. Центр КТ
удалён от центра частицы на
\begin{equation}
  R(g) = s + r_{\rm QD} + g,
  \qquad
  s = \begin{cases} c, & \text{КТ у вершины},\\ a, & \text{КТ сбоку}, \end{cases}
  \label{eq:center-distance}
\end{equation}
причём во втором случае центр КТ лежит в экваториальной плоскости.

Внешнее поле линейно поляризовано. Комбинация «размещение КТ $\times$
ориентация поля» даёт пять геометрически различных каналов
(табл.~\ref{tab:channels}); они различаются не только величиной отклика, но и
тем, нормально ли поле к поверхности МНЧ в точке, где стоит КТ.
В каждом канале рассматривается скалярный переход с дипольным моментом
вдоль выбранной компоненты поля. Пять каналов не являются готовым
усреднением по ориентациям кристаллических осей и переходных диполей
реальных коллоидных КТ.

\begin{table}[t]
  \centering
  \caption{Пять каналов возбуждения. «Ориентация» указана относительно оси
  симметрии сфероида, «к поверхности» — относительно локальной нормали в точке
  расположения КТ.}
  \label{tab:channels}
  \small
  \begin{tabular}{llll}
    \toprule
    Обозначение & Размещение КТ & Поле к оси & Поле к поверхности \\
    \midrule
    \texttt{axis\_long}              & вершина & вдоль   & нормально \\
    \texttt{axis\_trans}             & вершина & поперёк & касательно \\
    \texttt{side\_long}              & сбоку   & вдоль   & касательно \\
    \texttt{side\_trans\_radial}     & сбоку   & поперёк, радиально  & нормально \\
    \texttt{side\_trans\_tangential} & сбоку   & поперёк, касательно & касательно \\
    \bottomrule
  \end{tabular}
\end{table}

Сетка зазоров $g=0.5,\,0.75,\,1,\,1.5,\,2,\,3,\,5,\,7,\,10,\,12$ нм,
пять каналов и критерии отбора являются условиями настоящего исследования,
а не экспериментальными данными из~\cite{shah2013}.

\subsection{Квантовая точка}
\label{sec:qd}

КТ описывается как двухуровневая система с энергией перехода
$\hbar\omega_0=\SI{2.042}{\electronvolt}$ и эффективным дипольным моментом
перехода $d=\SI{13.9}{\debye}$. Феноменологические скорости релаксации:
населённости $\hbar\gamma_1=\SI{268}{\nano\electronvolt}$
($T_1\simeq\SI{2.46}{\nano\second}$) и чистой дефазировки
$\hbar\gamma_\varphi=\SI{1.270}{\milli\electronvolt}$.
Энергия перехода, диполь и обе скорости взяты из раздела II.D
работы~\cite{shah2013}. В той работе $d$ получен подгонкой оптического
отклика, а собственная скорость излучательного распада оценена по формуле
$\gamma_{\rm rad}=n_m\omega_0^3d^2/(3\pi\varepsilon_0\hbar c_0^3)$,
где $c_0$ --- скорость света в вакууме. При округлённых входах эта формула
даёт $\hbar\gamma_{\rm rad}\simeq\SI{267.3}{\nano\electronvolt}$;
в расчёте сохранено опубликованное значение \SI{268}{\nano\electronvolt}.
Полная скорость расфазировки равна
\begin{equation}
  \Gamma_2 = \gamma_\varphi + \tfrac{1}{2}\gamma_1,
  \qquad
  \hbar\Gamma_2 = \SI{1.270134}{\milli\electronvolt}.
  \label{eq:gamma2}
\end{equation}
Здесь коэффициент $\gamma_2$ при втором диссипаторе в уравнении (3)
работы~\cite{shah2013} прочитан как чистая дефазировка $\gamma_\varphi$:
распад населённости добавляет к затуханию когерентности $\gamma_1/2$.
В полуклассических уравнениях (5) той работы используется обозначение
$\gamma_2$ непосредственно для затухания когерентности. Поэтому конвенция
настоящего расчёта указана явно; численная разница здесь составляет
лишь \SI{0.000134}{\milli\electronvolt}. Получается
$T_2=1/\Gamma_2\simeq\SI{518}{\femto\second}$.
Выбранная дефазировка относится к принятому в~\cite{shah2013}
температурному масштабу $50$--$100$ К. Цитируемое там исследование
Besombes и соавторов~\cite{besombes2001} подтверждает роль акустических
фононов, но относится к CdTe/ZnTe и показывает также нелоренцевы боковые
полосы; оно не является измерением постоянной скорости для нашей CdSe-КТ.

В уравнениях динамики и квазистатической связи ниже используются атомные
единицы ($\hbar=e=m_e=4\pi\varepsilon_0=1$); размерные входы и результаты
приведены в эВ, нм, фс и Дж/см$^2$. Формула нормировки
флюенса~\eqref{eq:fluence} записана в СИ.
Уравнения Блоха записаны для вещественного поля, \emph{без} приближения
вращающейся волны. В переменных $W=\rho_{ee}-\rho_{gg}$ и вещественных
компонентах когерентности $P=2\operatorname{Re}\rho_{ge}$,
$Q=-2\operatorname{Im}\rho_{ge}$
\begin{align}
  \dot W &= \Omega(t)\,Q - \gamma_1\,(W+1), \label{eq:bloch-w}\\
  \dot Q &= -\omega_0 P - \Omega(t)\,W - \Gamma_2 Q, \label{eq:bloch-q}\\
  \dot P &= \omega_0 Q - \Gamma_2 P, \label{eq:bloch-p}
\end{align}
где мгновенная частота Раби определена через \emph{полное} поле в месте КТ,
\begin{equation}
  \Omega(t) = 2 d\,E_{\rm eff}(t),
  \qquad
  E_{\rm eff}(t) = f_{\rm loc}\bigl[E_{\rm inc}(t) + E_{\rm MNP}(t)\bigr].
  \label{eq:rabi}
\end{equation}
Дипольный момент КТ равен $\mu_d = f_{\rm loc}\,d\,P$, населённость верхнего
уровня $\rho_{ee}=(1+W)/2$. Фактор локального поля внутри КТ принят
$f_{\rm loc}=1$: дипольный момент задан во «внешней» конвенции, то есть уже
включает в себя диэлектрический отклик самой точки.
Начальное состояние --- $W=-1$, $P=Q=0$, все вспомогательные
амплитуды и скорости металла равны нулю. Хранящееся во входах
$\varepsilon_{\rm QD}=6$ не вносит дополнительного экранирования
при этой конвенции и не заимствовано из~\cite{shah2013}.

Для выбранной несущей длительность импульса
(раздел~\ref{sec:pulse}) составляет около десяти оптических периодов.
Решение для вещественного поля позволяет непосредственно получать
временной диполь и его спектр; отдельная оценка ошибки приближения
вращающейся волны в данной работе не выполняется.

Параметры $\gamma_1$ и $\Gamma_2$ одинаковы для всех геометрий.
Когерентное обратное действие металла включено через $E_{\rm MNP}$,
однако геометрически зависимые скорости спонтанного излучения
и безызлучательного распада некогерентно возбуждённой КТ отдельно
не вычисляются. Например, при нулевом поле и исчезнувших когерентности
и модальных амплитудах уравнения дают $\dot\rho_{ee}=-\gamma_1\rho_{ee}$
независимо от зазора. Это существенное отличие от полуклассической модели
с поправкой на эффект Парселла, обсуждаемой в~\cite{shah2013}.

\subsection{Дисперсия золота}
\label{sec:material}

Отклик МНЧ строится из табличных оптических констант золота
(Johnson--Christy~\cite{johnson1972}); между табличными энергиями линейно
интерполируются $n$ и $\kappa$, затем вычисляется
$\varepsilon_p=(n+i\kappa)^2$. Экстраполяция за таблицу не используется.
В квазистатике безразмерная поляризуемость сфероида вдоль
выбранной оси есть
\begin{equation}
  \alpha_L(\omega) =
  \frac{\varepsilon_p(\omega)-\varepsilon_m}
       {\varepsilon_m + L\,[\varepsilon_p(\omega)-\varepsilon_m]},
  \label{eq:alpha-L}
\end{equation}
где $L$ — фактор деполяризации соответствующей моды.

Для перехода во временную область нужна причинная и пассивная рациональная
модель отклика. Поэтому $\alpha(\omega)$ приближается суммой лоренцевых членов
\begin{equation}
  \alpha(\omega) = \alpha_\infty
  + \sum_{k=1}^{N} \frac{f_k}{\omega_k^2 - \omega^2 - i\gamma_k\omega},
  \qquad f_k \ge 0,\quad \gamma_k > 0,\quad \omega_k>0,
  \label{eq:lorentz-fit}
\end{equation}
причём $f_k\ge 0$ обеспечивает пассивность, а $\gamma_k>0$ — расположение
полюсов в нижней полуплоскости, то есть причинность при соглашении
$e^{-i\omega t}$. Подгонка выполняется в окне
$0.8$--$\SI{3.5}{\electronvolt}$, содержащем $22$ исходные табличные точки,
на сгущённой сетке их интерполяции с дополнительным контролем вблизи
\SI{2.042}{\electronvolt}. Сгущение не добавляет независимых измерений.
Принятое правило числа параметров $2\cdot22\ge (1+3N)+4$
ограничивает $N\le13$; это техническое условие переопределённости,
а не доказательство уникальности подгонки или число физических резонансов.
Продольная и поперечная поляризуемости подгоняются отдельно.
Величины $\omega_k,\gamma_k,f_k$ являются результатом этой процедуры;
плазмонная ширина \SI{150}{\milli\electronvolt} и диполь димерной
антенны из~\cite{shah2013} не служат их входами.

Число полюсов $N$ — первая из двух осей сравнения настоящей работы. Сравниваются
одноосцилляторная модель $N=1$, обычная в литературе по гибридным системам
КТ--МНЧ, и многополюсная $N=12$. Качество подгонки контролируется одновременно
по $\alpha$ и по $1/\alpha$, поскольку именно обратная величина входит в
знаменатели связанных уравнений.

\subsection{Связь КТ--МНЧ: три скалярных коэффициента отклика}
\label{sec:coupling}

Взаимодействие КТ и МНЧ в квазистатике полностью задаётся тремя
частотно-зависимыми величинами $A$, $B$, $K$:
\begin{equation}
  p_M=A(\omega)E_{\rm inc}+B(\omega)\mu_d,\qquad
  E_{\rm MNP}=B(\omega)E_{\rm inc}+K(\omega)\mu_d.
  \label{eq:ports}
\end{equation}
Здесь $A$ задаёт дипольный отклик МНЧ на однородное поле, $B$ ---
перекрёстную связь, $K$ --- поле, возвращаемое МНЧ к источнику-КТ.
Одинаковый $B$ в обоих выражениях отражает взаимность. Для пяти
симметричных конфигураций это скалярные проекции поля и диполя;
произвольная трёхмерная ориентация в данном расчёте не реализована.

\paragraph{Дипольное приближение (DD).} МНЧ заменяется одним точечным диполем
в её центре. Тогда все три величины выражаются через одну:
\begin{equation}
  A = C\,\alpha(\omega), \qquad B = A J, \qquad K = A J^2,
  \label{eq:dd}
\end{equation}
где $C=\varepsilon_m a^2 c/3$, а $J=G/(\varepsilon_m R^3)$,
$G=3(\mathbf e\cdot\widehat{\mathbf R})^2-1$; $\mathbf e$ --- направление
поля и диполя КТ, $\widehat{\mathbf R}$ --- направление между центрами.
Таким образом, $G=2$ для \texttt{axis\_long} и
\texttt{side\_trans\_radial}, $G=-1$ для остальных трёх каналов.
Продольность относительно оси сфероида сама по себе не определяет $G$.
Существенно, что в~\eqref{eq:dd} $K$ жёстко связано с
$A$: пространственная структура ближнего поля не содержит независимой
информации.

\paragraph{Полное квазистатическое решение (FQS).} Уравнение Лапласа решается
разделением переменных в вытянутых сфероидальных координатах
$(\xi,\eta,\phi)$. Пространственную гармонику обозначим
$\nu=(n,m,s)$, где $n\ge1$ --- степень, $0\le m\le n$ --- азимутальный
порядок, $s$ выбирает вещественный сектор $\cos(m\phi)$ или $\sin(m\phi)$.
Отрицательные $m$ отдельно не суммируются. Точечная КТ возбуждает
разрешённые симметрией гармоники множества $\mathcal M_h$ для канала $h$:
\begin{center}\small
\begin{tabular}{lll}
  \toprule
  Канал & Азимутальные порядки & Сектор \\
  \midrule
  \texttt{axis\_long} & $m=0$ & $\cos$ \\
  \texttt{axis\_trans} & $m=1$ & $\cos$ \\
  \texttt{side\_long} & $n+m$ нечётно & $\cos$ \\
  \texttt{side\_trans\_radial} & $n+m$ чётно & $\cos$ \\
  \texttt{side\_trans\_tangential} & $n+m$ чётно, $m\ge1$ & $\sin$ \\
  \bottomrule
\end{tabular}
\end{center}
Для бокового размещения центр КТ выбран на оси $x$; поля трёх каналов
направлены соответственно вдоль $z$, $x$, $y$. Однократная сумма по $n$
достаточна только для осевых каналов с фиксированным $m$.
Реакционное поле имеет вид
\begin{equation}
  K_h(\omega)=\lim_{n_{\max}\to\infty}
  \sum_{\substack{\nu=(n,m,s)\in\mathcal M_h\\n\le n_{\max}}}
       w^{(h)}_\nu\,\chi_{nm}(\omega),
  \qquad
  \chi_{nm}=\frac{\varepsilon_p-\varepsilon_m}
  {\varepsilon_m+L_{nm}(\varepsilon_p-\varepsilon_m)},
  \label{eq:K-series}
\end{equation}
где $w^{(h)}_\nu\ge0$ зависят от формы, положения и направления диполя,
а также от нормировки через $\varepsilon_m$. В расчёте ряд обрезается
при конечном $n_{\max}$ и проверяется его сходимость. Для поверхности
$\xi_0=c/\sqrt{c^2-a^2}$ факторы деполяризации определяются радиальными
присоединёнными функциями Лежандра:
\begin{equation}
  L_{nm}=\left.
  \frac{(P_n^m)'/P_n^m}{(P_n^m)'/P_n^m-(Q_n^m)'/Q_n^m}
  \right|_{\xi=\xi_0}.
  \label{eq:modal-L}
\end{equation}
Штрих здесь означает производную по $\xi$, а $P_n^m,Q_n^m$
не являются переменными Блоха. Сферический предел реализован отдельно,
$L_{nm}\to n/(2n+1)$.

Однородное поле и полный диполь МНЧ связаны только с одной светлой
гармоникой $b=(1,m_b,s_b)$ данного скалярного канала: $m_b=0$ при
поле вдоль $z$, $m_b=1$ при поперечном поле. Если $H=\chi_b$ и $J_b$
--- коэффициент связи этой гармоники с КТ, то
\begin{equation}
  A=CH,\qquad B=CJ_bH,\qquad K_b=CJ_b^2H=B^2/A.
  \label{eq:fqs-bright}
\end{equation}
При конечном расстоянии от вытянутого сфероида $J_b$ отличается от
точечно-дипольного $J$; они совпадают в сферическом или дальнем пределе.
Поэтому удержание только $n=1$ в FQS ещё не означает перехода к DD.

Ключевой элемент реализации — восприимчивости всех порядков выражаются через
\emph{одну и ту же} подгонку~\eqref{eq:lorentz-fit}. Исключая общую
проницаемость $\varepsilon_p$ из~\eqref{eq:alpha-L}, записанного для факторов
деполяризации $L_{nm}$ и $L_b$, получаем точное соотношение
\begin{equation}
  \chi_{nm} = \frac{H}{1 + (L_{nm} - L_b)\,H},
  \label{eq:modal-transform}
\end{equation}
где $H$ — подогнанная восприимчивость светлой моды. Соотношение реализуется как
причинная обратная связь, а не как независимые подгонки $A$, $B$ и $K$. Поэтому
светлые каналы остаются взаимными, а равенство $K_b = B^2/A$ выполняется
тождественно, по построению. Это устраняет риск того, что различие DD и FQS
окажется артефактом двух разных подгонок одного материала: в обеих моделях
материальный коэффициент буквально один и тот же, что проверяется численно
(раздел~\ref{sec:verification}).

Во временной области каждой удержанной пространственной моде $\nu$
соответствуют переменные $(q_{\nu k},\dot q_{\nu k})$ материальных
осцилляторов $k=1,\ldots,N$. При $\Delta L_\nu=L_\nu-L_b$ реализация имеет вид
\begin{align}
  v_\nu&=\frac{u_\nu-\Delta L_\nu\sum_kq_{\nu k}}
                   {1+\Delta L_\nu\alpha_\infty},\qquad
  h_\nu=\alpha_\infty v_\nu+\sum_kq_{\nu k},\label{eq:ade-feedback}\\
  \ddot q_{\nu k}+\gamma_k\dot q_{\nu k}+\omega_k^2q_{\nu k}
       &=f_kv_\nu.\label{eq:ade}
\end{align}
Входы равны $u_b=E_{\rm inc}+J_b\mu_d$ и $u_{\nu\ne b}=\mu_d$;
выходы --- $p_M=Ch_b$ и
$E_{\rm MNP}=CJ_bh_b+\sum_{\nu\ne b}w_\nu h_\nu$.
В DD остаётся одна светлая реализация с $J_b$, заменённым на $J$.
Реакция совокупности тёмных мод представляется меньшим числом эффективных
мод с положительными весами; её частотный отклик проверяется на отдельной
контрольной сетке. Например, для ближнего торцевого продольного сценария
$159$ тёмных гармоник заменены $16$ эффективными модами;
максимальная ошибка отклика, нормированная на его максимум в полосе,
равна $9.38\cdot10^{-7}$. Эффективные моды редукции не имеют определённых
индексов $(n,m)$. Увеличение $n_{\max}$ сохраняет эту редукцию и само по себе
не является сравнением нелинейных порогов с полностью нередуцированной
системой; такого отдельного сравнения в сохранённом прогоне нет.

\subsection{Возбуждающий импульс}
\label{sec:pulse}

Падающее поле — гауссов импульс с несущей $\omega_L$:
\begin{equation}
  E_{\rm inc}(t) = E_0 \exp\!\left(-\frac{t^2}{2\sigma^2}\right)\cos(\omega_L t),
  \qquad
  \sigma = \frac{\tau^{(I)}_{\rm FWHM}}{2\sqrt{\ln 2}},
  \label{eq:pulse}
\end{equation}
где $\tau^{(I)}_{\rm FWHM}=\SI{20}{\femto\second}$, как в разделе III.B
работы~\cite{shah2013}, --- полная ширина по
интенсивности (не по полю и не $\sigma$). Соответствующая спектральная ширина
по интенсивности равна $\SI{91.25}{\milli\electronvolt}$ — величина, которая
ниже окажется определяющей при сравнении наблюдаемых.

Амплитуда задаётся через падающий флюенс, вычисляемый точно для вещественного
гауссова импульса:
\begin{equation}
  \mathcal{F} = n_m\varepsilon_0 c_0 \int E_{\rm inc}^2(t)\,dt
  = \tfrac{1}{2}\sqrt{\pi}\,n_m\varepsilon_0 c_0\,\sigma E_0^2
    \left[1 + e^{-(\omega_L\sigma)^2}\right],
  \qquad n_m=\sqrt{\varepsilon_m}.
  \label{eq:fluence}
\end{equation}
Перекрёстный член в квадратных скобках сохранён, поэтому~\eqref{eq:fluence}
остаётся точным и для малоцикловых импульсов.
Соответствующая пиковая интенсивность усреднённой по оптическому циклу
огибающей равна $I_{\rm peak}=n_m\varepsilon_0c_0E_0^2/2$.

\subsection{Наблюдаемые величины}
\label{sec:observables}

\paragraph{Слабополевой спектр.} Линейный отклик КТ
\begin{equation}
  S(E) = \left|\frac{p_{\rm QD}(E)}{E_{\rm inc}(E)}\right|^2
  \label{eq:S}
\end{equation}
вычисляется алгебраически из линейного отклика связанной системы. По $S(E)$
определяются положение $E_{\rm res}$, ширина $\Gamma_{\rm eff}$ (по интервалу
половины проминенции) и сдвиг $\Delta E_{\rm res}$ относительно изолированной
КТ, чья ширина $\Gamma_0=\SI{2.54}{\milli\electronvolt}$ служит единицей.
Если $\beta(E)$ -- линейная поляризуемость изолированной КТ, то
\begin{equation}
  \frac{p_{\rm QD}}{E_{\rm inc}}
  = \frac{\beta(1+B)}{1-\beta K},
  \qquad
  \beta=\frac{2f_{\rm loc}^{\,2}d^2\omega_0}
  {\omega_0^2+(\Gamma_2-i\omega)^2}.
  \label{eq:linear-qd}
\end{equation}
Числитель описывает интерференцию падающего и наведённого полей,
знаменатель -- обратное действие МНЧ. Величина $S$ характеризует
когерентный дипольный отклик, а не квантовый выход или непосредственно
регистрируемую интенсивность люминесценции.
В карте по зазорам (рис.~\ref{fig:gap-spectral}) оба ядра используют
непосредственно интерполированный табличный материал, без лоренцевой
подгонки. Её влияние проверяется отдельно на рис.~\ref{fig:spectrum-material};
во временных расчётах DD и FQS используют общую подгонку с $N=12$.

Окно анализа признака равно $\mathcal W=[1.872,\,2.212]$ эВ,
энергетическая сетка покрывает $[1.85,\,2.234]$ эВ и содержит $4001$ узел.
Выбирается ближайший к \SI{2.042}{\electronvolt} пик с проминенцией
не менее $10\,\%$ от наибольшей в окне; его положение уточняется
параболой по трём соседним узлам. Проминенция --- превышение вершины
над более высоким из двух окружающих минимумов. Ширина измеряется
на половине этого превышения с линейной интерполяцией пересечений.
При обрезанном окном пике или конкурирующем пике с проминенцией
не менее половины выбранной единственная ширина не присваивается.

\paragraph{Усиление возбуждения.} Различаются две величины, отвечающие двум
разным экспериментам:
\begin{equation}
  G_{\rm exc} = \frac{S_{\rm hyb}(E_L)}{S_{\rm iso}(E_L)},
  \qquad
  \tilde G_{\rm exc} = \frac{\max_{E\in\mathcal W} S_{\rm hyb}(E)}
                            {\max_{E\in\mathcal W} S_{\rm iso}(E)}.
  \label{eq:gain}
\end{equation}
Первая отвечает узкополосному зондированию на \emph{фиксированной} частоте
$E_L=\hbar\omega_0$, вторая — перестройке лазера вслед за смещённым гибридным
резонансом внутри окна $\mathcal W$. Максимумы берутся на сохранённой
сетке; это операционная характеристика конечного окна.

\paragraph{Порог нелинейного возбуждения.} Основная наблюдаемая работы —
наименьший падающий флюенс $\mathcal{F}_\eta$, при котором на первой
возрастающей ветви Раби достигается
\begin{equation}
  \rho_{ee}\bigl(t_{\rm read};\,\mathcal{F}_\eta\bigr) = \eta = 0.5 .
  \label{eq:threshold}
\end{equation}
Момент считывания $t_{\rm read}=\SI{3600}{\femto\second}$ относительно
центра импульса общий для всех моделей, зазоров и каналов. К этому моменту
контролируется затухание когерентного отклика; собственная релаксация
населённости с принятой $\gamma_1$ остаётся малой.
Если цель не достигнута до верхнего флюенса, результат является нижней
границей порога. Если первая возрастающая ветвь закончилась ниже $\eta$,
фиксируется недостижение цели на этой ветви. Ни один из этих случаев
не означает невозможности возбуждения при произвольном импульсе.

\paragraph{Меры расхождения моделей.} Для переноса сравнения DD/FQS на сами
наблюдаемые вводятся относительные расхождения спектра $\delta_{\rm spec}(g)$ и
порога $\delta_{\mathcal F}(g)$; граница применимости DD определяется как
наименьший зазор, начиная с которого расхождение не превышает объявленного
допуска \emph{при всех} больших допустимых зазорах.
Граница относится к дискретной сетке и считается установленной,
только если после неё имеется хотя бы ещё один допустимый узел.

\subsection{Численная процедура и контроль сходимости}
\label{sec:numerics}

Система~\eqref{eq:bloch-w}--\eqref{eq:bloch-p} вместе с модальными
переменными интегрируется методом DOP853 с относительным допуском $10^{-7}$ и
абсолютным $10^{-9}$; шаг дополнительно ограничен сверху по числу точек на
период наибыстрейшей возбуждённой компоненты (не менее восьми точек).
В основной карте порог~\eqref{eq:threshold} находится последовательным
поиском первого пересечения по $\sqrt{\mathcal F}$ и методом Брента;
относительный допуск корня равен $10^{-5}$. Диапазон флюенсов составляет
$5\cdot10^{-10}$--$2\cdot10^{-4}$ Дж/см$^2$. Контроль изолированной КТ
использует $65$ узлов, а материальная серия рис.~\ref{fig:fluence-material}
--- $129$ узлов и интерполяцию порога по $\sqrt{\mathcal F}$.

Контроль включает уточнение пространственного порядка, уменьшение
максимального шага по времени, сгущение сеток по энергии и флюенсу
и увеличение числа материальных полюсов. Первоначальный порядок $80$
не прошёл проверку бокового продольного канала при
$g=\SI{0.5}{\nano\meter}$. Сохранённый основной расчёт использует
$n_{\max}=160$, а проверка выбранных конфигураций --- $320$.
Внутренний контроль ряда дополнительно сравнивает порядок с его половиной;
для основного расчёта допуск по реакционному ядру равен $2\cdot10^{-5}$.
Спектральные и пороговые проверки относятся
к соответствующим этапам расчёта; устойчивость рекомендации проверяется
на выбранных конфигурациях. Изменение результата сравнивается
с объявленным допуском. Помимо
этого контролируются пассивность реализации, положительность матрицы
плотности, доля спектра импульса внутри окна подгонки и затухание отклика к
моменту считывания. Обязательные проверки каждого основного этапа
останавливают его при нарушении допуска. Проверки чувствительности
сохраняют также отрицательные результаты; сканирование несущей
не включено в условие принятия ранжирования при номинальной частоте.
Для вариаций формы отдельно объявлены более мягкие допуски:
NRMS модального преобразования $0.05$, максимальная относительная
ошибка $0.12$, пространственная ошибка $10^{-4}$.
Они не заменяют допуски основного прогона.

\subsection{Верификация реализации}
\label{sec:verification}

Внутренние проверки сходимости не способны обнаружить ошибку в самих формулах
ядра, поэтому аналитические ядра сверяются с \emph{независимым} решателем
методом граничных элементов (декартовы поверхностные заряды), не использующим
ни сфероидальных гармоник, ни факторов деполяризации, ни аналитического
квазистатического отклика. Сверка охватывает все пять каналов
табл.~\ref{tab:channels} и сферический предел $c=a$; используется вложенная
последовательность сеток с экстраполяцией Ричардсона и объявленной
неопределённостью. Принимаются относительные расхождения не более $0.5\,\%$ для
$A$ и $B$ и $1\,\%$ для $K$, причём одновременно требуется попадание в
$3\sigma$ оценённой неопределённости метода граничных элементов.
Эта сверка не удостоверяет такую же точность при самых малых зазорах:
в дополнительной проверке при расстоянии от центра КТ до поверхности
МНЧ $0.25a$ собственная неопределённость гранично-элементного расчёта
по $K$ достигает $20$--$29\,\%$. Высокие пространственные порядки
на малых расстояниях дополнительно проверены по независимому
аналитическому ряду для сферического предела. Такая проверка
контролирует реализацию ряда, но не заменяет точный независимый
эталон для вытянутого сфероида при каждом зазоре.

Отдельно проверяется, что материальный коэффициент в DD и FQS совпадает с
машинной точностью, то есть расхождение этих двух моделей при общем
материале имеет пространственную природу. В проекте также предусмотрены отдельные сценарии
сопоставления с литературными моделями, включая~\cite{shah2013}; их геометрии
и параметры не подменяют сценарий с одной МНЧ.

\subsection{Область применимости}
\label{sec:applicability}

Модель локальна, квазистатична и трактует КТ как точечный диполь.
Помимо принятого описания релаксации, существенны следующие ограничения.

\emph{Квазистатика.} Запаздывание не учитывается; допустимыми считаются только
конфигурации с $k_m R \le 0.5$ и $k_mc\le0.3$, где
$k_m=\sqrt{\varepsilon_m}\,E/(\hbar c_0)$, а $c$ --- большая полуось МНЧ.
Для отбора используется верхняя контрольная энергия
\SI{2.234}{\electronvolt}, соответствующая верхней границе окна
рассчитанных спектров и превышающая несущую.
Это объявленный критерий отбора, а не доказанная граница точности
решения уравнений Максвелла.
Порядок отбрасываемого эффекта оценён отдельно: диагностическая поправка
Мейера--Вокауна~\cite{meier1983} даёт
множитель $1.29$ для $|\alpha|^2$ на несущей и сдвиг плазмонного резонанса
$-\SI{21}{\milli\electronvolt}$.

\emph{Точечность КТ.} Поле не усредняется по объёму КТ. Для радиальной
зависимости дипольного поля $E_{\rm dip}\propto R^{-3}$ индикатор его
изменения на размере точки равен $3r_{\rm QD}/R$, то есть $0.32$--$0.35$
в торцевой геометрии при $g=0.5$--$\SI{1.5}{\nano\meter}$.
Это оценка масштаба неоднородности дипольной компоненты, а не строгая
граница для полного поля $E_{\rm inc}+E_{\rm MNP}$, которое может
интерферировать. Малость поправки на конечный размер КТ не удостоверена;
вычисление поля в центре остаётся самостоятельным приближением.

\emph{Локальный отклик металла.} Нелокальное экранирование, туннелирование и
перенос заряда отсутствуют, поверхностное затухание оценено, но не включено
(множитель $0.83$--$0.95$ для $|\alpha|^2$). В этой отдельной оценке
используется $\Delta\gamma=A_s v_F/L_{\rm eff}$,
$L_{\rm eff}=4V/S_{\rm surf}$, где $V$ и $S_{\rm surf}$ --- объём
и площадь поверхности металла. Приняты $v_F=1.40\cdot10^6$ м/с
и $A_s\in[0.3,1]$; значение $A_s=0.3$ получено для золотых наностержней
в~\cite{novo2006}, а верхний предел $1$ здесь является выбором диапазона
оценки. Для выделения друдевской части заданы
$\hbar\omega_p=\SI{9}{\electronvolt}$ и
$\hbar\gamma_D=\SI{0.07}{\electronvolt}$: это значения того же масштаба,
что экспериментальные $\hbar\omega_p=\SI{8.45}{\electronvolt}$ и
$\tau_D=14\pm3$ фс в~\cite{olmon2012}, но не их точное заимствование.
Все эти параметры относятся только к диагностике отброшенных эффектов;
основные кривые используют таблицу Johnson--Christy без данной поправки.
При субнанометровых зазорах
особенно важно проверить допустимость локального отклика и отсутствие
переноса заряда; универсальная граница по $g$ без описания оболочек
и электронного барьера не задаётся. Поверхностное затухание зависит также
от размеров МНЧ и не исчезает при увеличении зазора.

Сравнения DD/FQS и $N=1$/$N=12$ выполняются при одинаковых остальных
параметрах, что позволяет выделить ошибку каждого упрощения внутри
принятой модели. Однако неучтённые физические эффекты не обязаны
сокращаться в отношениях откликов: они могут менять резонансы и
геометрическую зависимость связи. Численная сходимость не заменяет
проверки этих систематических погрешностей.
